\chapter{通往主动推断的低路}

\section{引言}

本章从亥姆霍兹式, 或者也可以说康德式的观点出发来介绍主动推断: 知觉乃是``无意识的推断''(Helmholtz 1867)。同时, 本章也会讨论在贝叶斯大脑假说之下近年发展出的相关思想。主动推断不仅把知觉视为推断问题, 还把行动、规划与学习一并处理为(贝叶斯)推断问题, 并为这些原本难以直接求解的问题给出一种有原则的(变分)近似方法。正是在这一点上, 主动推断既吸收又扩展了前述思想。

\section{作为推断的知觉}

把大脑看成一台``预测机器'', 或看成一个能够推断并预测外部世界状态的统计器官, 这一传统由来已久。这个想法可以追溯到``知觉是无意识推断''这一观念(Helmholtz 1866), 而在较近时期, 它又被重新表述为``贝叶斯大脑''假说(Doya 2007)。从这个视角看, 知觉并不是一种纯粹自下而上的转导过程, 即把感觉状态(例如来自视网膜的输入)直接变成关于外部事物的内部表征(例如神经元活动模式)。相反, 知觉是一种推断过程: 它把关于感觉最可能成因的自上而下先验信息, 与自下而上的感觉刺激结合起来。推断过程作用于对世界状态的概率表征, 并遵循贝叶斯法则; 贝叶斯法则规定了在感觉证据出现之后, 应如何以最优方式更新信念。知觉不是一个被动的、由外而内的信息提取过程, 不是简单从感觉上皮所受到的``外部''刺激中抽取信息; 它是一个建构性的、由内而外的过程, 在这个过程中, 感觉被用来证实或证伪关于其生成方式的假设(MacKay 1956; Gregory 1980; Yuille and Kersten 2006; Neisser 2014; A.~Clark 2015)。

而要执行贝叶斯推断, 我们需要一个生成模型, 它有时也被称为前向模型。生成模型是统计理论中的一种构造, 用来生成关于观测的预测。它可以写成观测 $y$ 与产生这些观测的世界隐状态 $x$ 的联合概率 $P(y, x)$。后者之所以被称为隐状态或潜在状态, 是因为它们无法被直接观测。这个联合概率可以分解为两部分。第一部分是先验 $P(x)$, 表示有机体在看到感觉数据之前, 对世界隐状态所拥有的知识。第二部分是似然 $P(y \mid x)$, 表示有机体关于观测如何由状态生成的知识。贝叶斯法则告诉我们如何把这两部分结合起来, 本质上就是在接收观测之后, 把先验概率 $P(x)$ 更新为关于隐状态的后验概率 $P(x \mid y)$。对于需要快速回顾基础概率论的读者, 框 2.1 给出了简要总结。

贝叶斯推断这一主题非常广泛, 它出现在统计学、机器学习与计算神经科学等多个学科中。对相关内容的完整处理超出了本书范围, 但希望深入理解的读者可以参考非常优秀的教材(Murphy 2012)。不过, 这一切都建立在一条简单规则之上。为了说明这条规则, 我们来看一个贝叶斯知觉推断的例子(图 2.1)。设想某人强烈相信自己面前的是一个苹果。这个信念对应于先验概率, 简称先验。这个先验包含分配给``苹果''假设的概率, 以及分配给替代假设的概率。在这个例子里, 替代假设是它不是苹果, 而是一只青蛙。在数值上, 先验概率分布给苹果赋值 $0.9$, 给青蛙赋值 $0.1$。注意, 由于我们假定只有两个合理且互斥的假设, 它们的概率必须加总为 $1$。与此同时, 这个人还拥有一个似然模型, 该模型给``青蛙会跳''赋予很高概率, 而``苹果会跳''的概率很低。这个似然给出了从两个隐状态(青蛙或苹果)到两个观测(跳或不跳)的概率映射。先验与似然合在一起, 就构成了该人的生成模型。

现在设想这个人观察到她面前这个``苹果--青蛙''跳了起来。贝叶斯法则告诉我们, 应当如何在考虑``跳跃''这一似然信息后, 由先验形成后验信念。这个规则写作:

\begin{equation}
P(x \mid y) = \frac{P(x)P(y \mid x)}{P(y)}
\tag{2.1}
\end{equation}

\paragraph{框 2.1 概率的求和规则与乘法规则}

概率推理依赖于两条关键规则: 概率的求和规则与乘法规则, 分别为

\[
\sum_x P(x) = 1
\]

\[
P(x)P(y \mid x) = P(x, y)
\]

求和规则表示, 所有可能事件 $x$ 的概率必须加总(或积分)为 $1$。乘法规则表示, 两个随机变量 $x$ 与 $y$ 的联合概率, 可以分解为其中一个变量的概率 $P(x)$ 与在该变量给定时另一个变量的条件概率 $P(y \mid x)$ 的乘积。所谓条件概率, 就是在我们知道一个变量(这里是 $x$)取何值的条件下, 另一个变量(这里是 $y$)发生的概率。

由这两条简单规则, 我们可以推出两个重要结果。第一个是边缘化操作, 第二个是贝叶斯法则。边缘化使我们能够从联合分布中得到其中一个变量的分布:

\[
\sum_x P(x, y) = \sum_x P(y)P(x \mid y) = P(y)\sum_x P(x \mid y) = P(y)
\]

其中, $P(y)$ 被称为边缘概率, 而这个操作则称为把 $x$ 边缘化掉。贝叶斯法则也可以直接从乘法规则得到:

\[
P(x)P(y \mid x) = P(x, y) = P(y)P(x \mid y)
\]

这使我们能够在先验与条件分布(似然), 以及相应的边缘分布与另一条件分布(后验)之间来回转换。简单说, 贝叶斯法则无非是在说: 两件事同时发生的概率, 等于在第二件事已知时第一件事发生的概率乘以第二件事的概率; 这也等于在第一件事已知时第二件事发生的概率乘以第一件事的概率。

在图 2.1 的似然模型下, 分配给青蛙的后验概率是 $0.9$, 分配给苹果的概率是 $0.1$。如框 2.1 所强调, 公式 2.1 的分母可以通过对分子进行边缘化来计算。借助苹果--青蛙这个例子, 我们也可以顺带展开说明两种不同的``惊讶''概念, 这两者在主动推断中都非常重要。

第一种, 我们直接称之为惊讶(surprise), 它是负对数证据, 而证据就是观测的边缘概率。在我们的例子里, 它就是在生成模型之下, 观察到某物跳起来这一事件的负对数概率。对贝叶斯视角而言, 惊讶是一个非常重要的量, 因为它衡量的是模型对其试图解释的数据拟合得有多差。直观地说, 我们可以算出在模型下观察到``跳跃''行为的概率。请记住, 这个模型给苹果赋予了很高的先验概率, 给青蛙赋予了较低的先验概率。因此, ``跳跃''的边缘概率是

\begin{equation}
\begin{aligned}
P(y = \text{jumps}) &= \sum_x P(x, y = \text{jumps}) \\
&= \sum_x P(x)P(y = \text{jumps} \mid x) \\
&= P(x = \text{frog})P(y = \text{jumps} \mid x = \text{frog}) \\
&\quad + P(x = \text{apple})P(y = \text{jumps} \mid x = \text{apple}) \\
&= 0.1 \times 0.81 + 0.9 \times 0.01 \\
&= 0.09
\end{aligned}
\tag{2.2}
\end{equation}

这意味着, 在这个模型下, 我们平均每 $100$ 次观测中只会大约有 $9$ 次看到跳跃行为。因此, 如果我们认同图 2.1 中的模型, 那么观察到跳跃就应当让我们感到惊讶。我们可以用惊讶 $\Im$ 来量化这一点:
\[
\Im(y = \text{jumps}) = - \ln P(y = \text{jumps}) = - \ln(0.09) = 2.4 \text{ nats}.
\]
这个数越大, 模型作为当前观测之解释就越差。这使我们能够相对于数据来比较不同模型。比如, 考虑一个替代模型: 我们先验上相信自己 $100\%$ 的时间都会看到青蛙。按照与公式 2.2 同样的步骤计算, 会得到大约 $0.2$ nats 的惊讶值。这个模型更适合这些数据, 因为这项观测在该模型下没有那么令人惊讶。基于证据(或惊讶)为模型打分的程序, 常被称为贝叶斯模型比较。对于更复杂的模型, 惊讶的形式可能不再这么简单。表 2.1 给出了一系列概率分布对应的惊讶形式(省略常数项), 除了我们例子中的范畴分布之外也一并列出。关键在于, 这让我们能够针对支持集不同于当前简单例子的概率分布来讨论惊讶。这一点很重要, 因为世界生成感觉数据的方式会随数据类型而变化。我们可能会因为看到一张本不期待见到的面孔而感到惊讶(范畴分布), 也可能因为室外温度比预想更低而感到惊讶(连续分布)。表 2.1 可以看作是我们在后续章节构建生成模型时可用的一组概率分布工具箱。更一般地说, 它说明惊讶是一个可以针对任意给定概率分布族进行评估的概念。

\paragraph{表 2.1 概率分布与惊讶}

\begin{center}
\begin{tabular}{lll}
\hline
分布 & 支持集 & 惊讶 $\Im$ \\
\hline
Gaussian & $x \in \mathbb{R}$ & $(x-\mu)^2$ \\
Multinomial & $x_i \in \{0,\ldots,N\},\ \sum_i x_i = N$ & $-\sum_i x_i \ln d_i$ \\
Dirichlet & $x_i \in (0,1),\ \sum_i x_i = 1$ & $\sum_i (1-\alpha_i)\ln x_i$ \\
Gamma & $x \in (0,\infty)$ & $bx + (1-a)\ln x$ \\
\hline
\end{tabular}
\end{center}

注: Multinomial 的特例包括 categorical($K>2, N=1$)、binomial($K=2, N>1$) 与 Bernoulli($K=2, N=1$) 分布; Dirichlet 的一个特例是 beta 分布($K=2$)。

第二种惊讶概念被称为贝叶斯惊讶(Bayesian surprise)。这个术语多少有些容易混淆, 因为它衡量的是在一次观测之后我们必须更新信念的程度。换言之, 贝叶斯惊讶量化的是先验分布与后验分布之间的差异。这就提出了一个问题: 我们应如何度量两个概率分布之间的不同? 信息论给出的一个答案是使用 Kullback--Leibler 散度(KL 散度)。它被定义为两个对数概率之差的平均值:

\begin{equation}
D_{\mathrm{KL}}[Q(x)\Vert P(x)] \triangleq E_Q(x)\bigl[\ln Q(x) - \ln P(x)\bigr]
\tag{2.3}
\end{equation}

这里的 $E$ 表示平均值(或期望), 具体见框 2.2。借助 KL 散度, 我们可以量化本例中的贝叶斯惊讶:

\begin{equation}
\begin{aligned}
D_{\mathrm{KL}}[P(x \mid y)\Vert P(x)]
&= P(x = \text{frog} \mid y = \text{jumps}) \bigl(\ln P(x = \text{frog} \mid y = \text{jumps}) - \ln P(x = \text{frog})\bigr) \\
&\quad + P(x = \text{apple} \mid y = \text{jumps}) \bigl(\ln P(x = \text{apple} \mid y = \text{jumps}) - \ln P(x = \text{apple})\bigr) \\
&= 0.9\bigl(\ln 0.9 - \ln 0.1\bigr) + 0.1\bigl(\ln 0.1 - \ln 0.9\bigr) \\
&\approx 1.8 \text{ nats}
\end{aligned}
\tag{2.4}
\end{equation}

这个量衡量的是信念更新的幅度, 而不是观测本身有多不可能。为了突出惊讶与贝叶斯惊讶之间的区别, 考虑这样一种情况: 我们先验上坚信自己总会看到苹果。在这种情况下, 贝叶斯惊讶将为零, 因为先验信念如此强烈, 以至于我们在看到观测之后完全不会更新它。然而, 惊讶却会非常大(4.6 nats), 因为苹果跳起来本身极不可能。

\paragraph{框 2.2 期望}

引用随机变量 $x$ 的期望很有用, 通常记为 $E[x]$。它表示该变量所有可能取值的加权平均, 权重即各取值对应的概率。对于离散随机变量(只能取可数多个值), 期望可写为加权求和:

\[
E[x] = \sum_x xP(x)
\]

例如, 若一个离散数值变量只能以相等概率 $1/2$ 取值 $1$ 和 $2$, 则
\[
E[x] = 1 \times \frac{1}{2} + 2 \times \frac{1}{2} = \frac{3}{2}.
\]

对于连续随机变量(可以取无限多个值), 求和需要替换为积分。期望也可以作用于随机变量的函数, 而不是变量本身。例如, 若有函数 $f(x)$, 且 $x$ 服从某个连续分布, 则

\[
E[f(x)] = \int f(x)p(x)\,dx.
\]

本书会一直使用这一记号, 其中函数 $f(x)$ 往往是对数概率, 或对数概率比。

请注意, 尽管我们是以一个非常简单的生成模型来展示贝叶斯推断, 但它实际上适用于任意复杂度的生成模型。在第 4 章中, 我们将重点介绍支撑主动推断大多数应用的两类生成模型。

\section{生物学推断与最优性}

有两个重要观点把上述推断框架与生物学和心理学中的知觉理论联系起来。第一, 前面讨论的推断程序需要自上而下过程与自下而上过程的相互作用: 前者编码来自先验的预测, 后者编码由似然所中介的感觉观测。正是这种自上而下与自下而上之间的互动, 使推断视角有别于那些只考虑自下而上过程的替代路径。此外, 这种互动也是现代生物学知觉理论的核心, 例如预测编码(将在第 4 章讨论), 它是这里所讨论的更一般的(贝叶斯)推断框架的一种特定算法性(或过程层面)实现。

第二, 贝叶斯推断是最优的。所谓最优性, 是相对于一个被优化(即最小化)的代价函数而言的; 对于贝叶斯推断, 这个代价函数被称为变分自由能, 它与惊讶密切相关。我们将在第 2.5 节回到这一点。通过显式考虑隐状态上的完整分布, 贝叶斯推断能够自然处理不确定性, 从而避免那些只考虑隐状态点估计(例如 $x$ 的均值)的方法所具有的局限。一种这样的替代方法是极大似然估计, 它只是简单选出最有可能生成当前数据的那个隐状态。问题在于, 这种估计忽略了隐状态的先验合理性, 也忽略了估计本身周围的不确定性。贝叶斯推断没有这些局限。

不过, 尽管惊讶可以用来客观评估一个模型是否适用, 我们仍需理解, 推断本身是主观的。推断结果未必在客观意义上准确, 也就是说, 有机体的信念未必真的对应现实, 原因至少有两点。第一, 生物体所拥有的计算资源与能量资源都是有限的, 这使得精确贝叶斯推断通常不可行。这就要求我们采用近似, 而近似会让精确贝叶斯最优性不再有保障。这些近似包括变分后验这一概念, 它建立在所谓平均场近似之上, 并在第 4 章处于核心位置。

第二个使最优性呈现出主观性质的原因在于, 有机体运行的基础是主体关于其观测如何产生的生成模型, 而这个模型未必对应于现实中真正生成这些观测的过程。这并不是说, 生成模型就必须与生成过程完全一致。实际上, 某些模型甚至可能比真实生成它们的过程更擅长给出更好的(例如更简单的)数据解释; 这一点可以通过相对惊讶来量化。一个很好的例子是错觉: 某人面对被一位调皮的心理物理学家精心设计过的视觉刺激时, 会为自己的视觉输入找到一个更简单的解释。随着新经验的积累, 生成模型本身也可能被优化, 而这种优化未必会收敛到真实的生成过程。

图 2.2 说明了这一点, 也展示了真实环境因果关系(即生成过程)与有机体对世界的生成模型之间的区别。在这个例子中, 生成过程处于一个对有机体不可达的真实状态 $x^\ast$。然而, 有机体与世界是相互耦合的, 因此 $x^\ast$ 会生成一个观测 $y$, 并被有机体感知。有机体随后可以利用这个观测 $y$ 与贝叶斯法则, 去推断生成模型中某个解释变量或隐状态的后验概率。图中我们同时把 $x^\ast$ 与 $x$ 都称作隐状态, 以强调两者都不可直接观测, 但它们又有细微差异: 前者属于生成过程, 后者属于有机体的生成模型。进一步说, $x^\ast$ 与 $x$ 并不一定生活在同一个状态空间里。外部世界中的隐状态可能会取到超出大脑可用解释空间的值; 反过来, 大脑的解释也可能包含一些外部世界并不存在的变量。例如, 前者可以是五维的, 后者是二维的; 或者一个是连续的, 另一个是范畴性的。行动 $u$ 则是依据生成模型下所做出的推断而产生的, 但一旦执行, 它就作为生成过程的一部分作用于世界。

生成模型与生成过程之间的区分, 对于把握心理学中关于推断最优性的主张非常重要。即便这些主张成立, 它们在贝叶斯视角下也总是以有机体可用资源为条件。这里所谓资源, 指的是其具体的生成模型, 以及受限的计算与记忆资源。

\section{作为推断的行动}

到目前为止的讨论, 与所有贝叶斯大脑理论是一致的。不过现在, 我们引入主动推断所提供的一个简单但根本性的推进: 把行动也视为推断。这个思想起始于这样一个观念: 贝叶斯推断会最小化惊讶(或者等价地, 最大化贝叶斯模型证据)。前面我们讨论的是, 当我们通过推断来计算惊讶, 并依照模型最小化惊讶的能力在模型之间做选择时, 会发生什么。然而, 惊讶不仅依赖于模型, 还依赖于数据本身。若我们通过作用于世界来改变数据的生成方式, 就可以挑选那些在我们的模型之下最不令人惊讶的数据, 从而保证模型适用。

一旦配备了产生行动的机制, 有机体就能够与其环境展开相互往复的交换, 见图 2.2。在动物身上, 这一机制表现为运动反射回路。本质上, 在每一个行动--知觉循环中, 环境向有机体发送一个观测; 有机体使用(对)贝叶斯推断的(近似)来推断其最可能的隐状态; 然后它生成一个行动, 并把行动发送给环境, 试图让环境变得没那么令人惊讶。环境执行这一行动, 产生新的观测, 再把它发送给有机体, 新的循环由此开始。这里采用顺序描述只是出于教学目的; 重要的是要看到, 这些并不真是离散步骤, 而是连续的动力学过程。

主动推断不只是认识到知觉与行动拥有相同的推断本性。它还主张, 知觉与行动是合作着实现同一个目标, 也就是共同优化同一个函数, 而不是像更常见的理论那样分别优化两个不同目标。在主动推断文献中, 这个共同目标曾被以多种非正式或正式方式描述为最小化惊讶、熵、不确定性、预测误差, 或(变分)自由能。这些术语彼此相关, 但其关系有时并不立刻清晰, 因而造成了一些混乱。更进一步, 它们往往出现在不同语境中; 例如, 预测误差最小化多出现在生物学语境里, 用于解释脑信号, 而变分自由能最小化则更多出现在机器学习语境中。

在接下来的两节里, 我们将澄清, 主动推断主体通过知觉与行动实际最小化的单一量是变分自由能。不过在某些条件下, 变分自由能可以还原为其他概念, 例如生成模型与世界之间的不匹配, 或我们所期望者与我们所观察者之间的差异(即预测误差)。第 2.5 节将正式引入变分自由能。为简明起见, 第 2.4 节先聚焦于知觉与行动如何最小化生成模型与世界之间的差异。

\section{最小化模型与世界之间的差异}

既然我们已经用贝叶斯推断刻画了知觉与行动, 现在就可以进一步追问: 推断的目标究竟是什么? 换句话说, 推断在优化什么? 在认知科学中, 人们通常假定不同认知功能, 如知觉与行动, 优化的是不同目标。例如, 可以认为知觉最大化重建准确性, 而行动选择最大化效用。与此相反, 主动推断的一个根本洞见是: 知觉与行动服务于完全相同的目标。作为第一近似, 这一共同目标可以表述为最小化模型与世界之间的差异。有时这一点会被操作化为预测误差最小化。

为了理解知觉与行动如何减少模型与世界之间的差异, 我们再来看那个期待看见苹果的人(图 2.3)。她会生成一个自上而下的视觉预测, 例如``会看到某个红色且不会跳动的东西''。这个视觉预测随后会与感觉输入进行比较, 比如感觉到``有某个东西跳了起来'', 而这种比较便产生了差异。

她有两种方式来消除这一差异。第一, 她可以改变自己对所见之物的判断, 承认那是青蛙而不是苹果, 以便让自己的心智适配世界, 从而消除差异。这对应于知觉。第二, 她也可以把视线转移到最近的苹果树上, 看到一个看起来确实很像苹果的东西。这样同样能消除最初的差异, 但方式不同: 它不是改变心智来适配世界, 而是改变世界(包括自己的注视方向)以及随后的感觉输入, 使之适配心中已有的预期。这就是另一种拟合方向, 也就是行动。

如果在苹果与青蛙的世界里, 改变目光方向听起来还没有改变想法那么有说服力, 那么考虑另一种情况: 某个人期待自己的体温维持在某个区间, 但通过中枢温度感受器却感知到高温。这样的状态会让人惊讶, 也带来了一个需要解决的显著差异。与前面的例子一样, 他有两条最小化差异的路径, 分别对应于知觉(改变想法)与行动(改变世界)。在这里, 仅仅``改变想法''显然并不适应, 而采取行动来降低体温(例如打开窗户)则是适应性的。

这也说明, 在主动推断中, 边缘概率或惊讶(例如对体温的惊讶)这一概念, 其意义超出了标准贝叶斯处理方式, 并吸收了诸如内稳态与异稳态设定点之类的思想。从技术上说, 主动推断主体拥有的模型会给它们偏好访问的状态, 或偏好获得的观测, 赋予高边缘概率。对鱼来说, 这意味着``处于水中''会具有很高的边缘似然。换句话说, 有机体会隐含地期待自己抽样到的观测落在其舒适区内, 例如处于生理边界允许的范围中。

总而言之, 我们讨论了在任意一个时间点上, 知觉与行动如何共同帮助我们最小化模型与世界之间的差异。我们究竟是调整信念还是调整数据, 取决于我们对这些信念有多大把握。在苹果的例子中, 相关信念带着足够高的不确定性, 因而更适合被更新而不是被付诸行动。相反, 在体温的例子中, 我们对核心体温拥有更高的确信, 因为它支撑着我们的存在本身。正因如此, 我们会更新世界, 使之符合我们的信念。不过, 在主动推断中, 知觉与行动之间的合作比上述表述还要紧密。为了理解这一点, 下一节将从较受限的``差异''(或预测误差)概念, 转向更一般的``变分自由能''概念; 后者才是主动推断真正最小化的量, 而预测误差只是它的一个特例。

\section{最小化变分自由能}

到目前为止, 我们是在一个试图最小化惊讶的贝叶斯框架下讨论知觉与行动的。然而, 在绝大多数情形中, 支撑知觉与行动的精确贝叶斯推断在计算上都是不可行的, 因为有两个量无法直接求出: 模型证据 $P(y)$ 与后验概率 $P(x \mid y)$。这可能有两个原因。第一, 对于复杂模型而言, 可能存在许多种都需要被边缘化掉的隐状态, 这会让问题在计算上变得不可处理。第二, 边缘化操作本身可能需要解析上无法求解的积分。主动推断诉诸一种对贝叶斯推断的变分近似, 从而使问题变得可处理。

第 4 章将详细展开变分推断的形式主义。这里我们只需知道, 执行变分贝叶斯推断意味着把两个难以计算的量, 即后验概率与(对数)模型证据, 替换为两个能够高效计算的近似量: 近似后验 $Q$ 与变分自由能 $F$。近似后验有时也被称为变分分布或识别分布。负的变分自由能在机器学习中也常被称作证据下界(evidence lower bound, ELBO)。

最重要的是, 贝叶斯推断问题因此被转化为一个优化问题: 最小化变分自由能 $F$。变分自由能这一量源于统计物理, 在主动推断中扮演根本性角色。在公式 2.5 中, 它记为 $F[Q, y]$, 因为它既是近似后验 $Q$ 的泛函(函数的函数), 也是数据 $y$ 的函数:

\begin{equation}
\begin{aligned}
F[Q, y]
&= -E_Q(x)[\ln P(y, x)] - H[Q(x)] \\
&= D_{\mathrm{KL}}[Q(x)\Vert P(x)] - E_Q(x)[\ln P(y \mid x)] \\
&= D_{\mathrm{KL}}[Q(x)\Vert P(x \mid y)] - \ln P(y)
\end{aligned}
\tag{2.5}
\end{equation}

这三行分别可以理解为:

\[
\text{Energy} - \text{Entropy}
\]

\[
\text{Complexity} - \text{Accuracy}
\]

\[
\text{Divergence} - \text{Evidence}
\]

乍看之下, 变分自由能似乎是一个抽象概念, 但只要把它分解为认知科学里更直观、更熟悉的量, 它的性质及其在主动推断中的作用就会变得清晰。对变分自由能的每一种表述, 都给我们提供了一个理解自由能最小化含义的有用直觉。这里我们先作简要勾勒, 因为它们在本书后半部分讨论实例时会变得重要。

公式 2.5 的第一行表明, 相对于 $Q$ 进行最小化时, 我们既要求与生成模型保持一致(能量项), 又要求保持较高的后验熵。后者意味着, 在没有数据或没有精确先验信念时(这些只会影响能量项), 我们应当对世界隐状态采取尽可能不确定的信念, 这与 Jaynes 的最大熵原理(Jaynes 1957)一致。简单说, 没有信息时就应当保持不确定。这里的``能量''一词继承自统计物理。具体说, 在玻尔兹曼分布下, 一个系统采取某种构型的平均对数概率, 与该构型所关联的能量成反比, 即系统从基线构型移动到该构型所需的能量。

第二行强调, 自由能最小化可以理解为为感觉数据寻找最佳解释, 而这个最佳解释必须既足够简单(复杂度最低), 又能够准确解释数据(参见奥卡姆剃刀)。复杂度与准确性之间的权衡在多个领域反复出现, 通常出现在数据分析中的模型比较语境里。在统计学中, 有时也会用其他近似模型证据的指标, 例如贝叶斯信息准则或赤池信息准则。当我们后面介绍如何在基于模型的数据分析中用自由能进行模型比较, 以及在结构学习与模型缩减的语境下讨论时, 复杂度--准确性权衡会变得很重要。从认知角度看, 推断出复杂度最小的解释同样重要, 因为可以假设, 为了容纳数据而更新我们已经知道的内容(即先验), 会带来认知代价(Ortega and Braun 2013; Z\'enon et al. 2019); 因此, 与先验偏离最小的解释会更可取。

从这个角度看, 复杂度代价其实就是贝叶斯惊讶。换句话说, ``我改变自己想法的幅度''可以用先验与后验之间的散度来量化。这意味着, 对感觉的每一个准确解释都要付出复杂度代价, 而这个代价衡量的正是贝叶斯信念更新的程度。因此, 变分自由能刻画的是准确性与复杂度之间的差值。

最后一行则把自由能表述为负对数证据的一个上界(见图 2.4)。图的左半部分说明, 自由能是负对数证据的上界, 其中上界与若能执行精确推断(而非变分推断)时本应获得的后验概率之间的差异, 正是 $Q$ 与真实后验之间的散度。图的右半部分表明, 随着这个散度减小, 自由能会逼近负对数证据(惊讶); 当近似后验 $Q$ 与精确后验 $P(x \mid y)$ 完全一致时, 它就与惊讶相等。这为知觉推断提供了形式化动机: 通过尽可能优化近似后验 $Q$, 我们就能降低自由能。

公式 2.5 的最后一行还表明, 知觉推断并不是最小化自由能的唯一方式。我们也可以通过行动来改变感觉数据, 从而改变对数证据项。从认知角度看, 这一分解很有意思, 因为最小化散度与最大化证据, 分别对应于知觉与行动这两个互补的子目标, 见图 2.5。注意, 如果把 $Q$ 用 $P(x \mid y)$ 替换, 那么上述所有表达式都会退化为对负对数证据的刻画, 也就是精确推断的情形。

总之, 主动推断就是通过知觉与行动最小化变分自由能。这种最小化使有机体能够让自己的生成模型去拟合它所抽样到的观测。这样的拟合, 一方面衡量知觉上的充分性(由散度项表达), 另一方面也衡量对外部状态的主动控制能力, 因为它使有机体能够把自己维持在由生成模型所定义的一组适宜且偏好的状态内。换一种表述, 可以借助自由能的散度/证据分解来理解。把负对数证据等同于惊讶, 再注意到散度最小只能到 $0$, 我们就能看到自由能是惊讶的一个上界, 因而它只可能大于或等于惊讶。当有机体通过知觉最小化散度时, 自由能就成为惊讶的一个近似; 当有机体进一步通过行动改变自己收集到的观测, 使其更接近先验预测时, 它就会最小化惊讶。

变分自由能具有回顾性, 因为它只依赖过去与当前的观测, 不依赖未来观测。尽管它可以基于过去数据帮助我们对未来作出推断, 但它并不直接支持基于预期未来数据的前瞻性推断形式。而这对于规划与决策至关重要。在这些场景下, 我们需要根据某个行动或行动序列(策略)预期会带来的未来观测, 去推断最好的行动。要做到这一点, 就必须在生成模型中补充``期望自由能''这一概念。

\section{期望自由能与作为推断的规划}

期望自由能把主动推断扩展到一种典型的前瞻性认知形式: 规划。规划一连串行动, 例如走出迷宫所需的一系列移动, 需要考虑我们预计将会收集到的未来观测。比如, 某些备选行动路线的后果可能包括: 右转之后看到一个死胡同, 或在连续三次左转之后看到出口。每一种可能的行动序列都被称为一项策略(policy)。这也凸显了主动推断中对行动与策略所作的一个重要区分。前者是直接作用于外部世界的东西, 后者则是一种关于行为方式的假设。其含义是, 主动推断把规划与决策视为一个推断``该做什么''的过程。这样一来, 规划就被明确纳入了贝叶斯推断的领域, 我们也必须像前面一样指明先验与似然(见第 2.1 节)。不过, 在这里, 可供选择的不再是青蛙与苹果, 而是行为策略(例如: 我更可能去看池塘还是看果树?)。

在这一节中, 我们先简要处理似然, 即执行某项策略所带来的后果, 然后再转向先验, 而期望自由能正是在这里登场。

与策略相关的结果并不是当下立刻可得的, 因为它们位于未来, 但我们可以把生成模型中的两个组成部分串联起来, 从而对其进行预测。第一部分是我们关于隐状态如何随策略而变化的信念。具体细节将在第 4 章讨论。此处我们用 $\tilde{x}$ 表示一段时间上的隐状态序列或轨迹, 并把轨迹条件化在生物所采取的策略 $\pi$ 之上。这意味着, 我们模型的动力学部分写作 $P(\tilde{x} \mid \pi)$。继续沿用苹果--青蛙的例子, 策略可以是去池塘还是去果园的决定, 而这会改变遇到青蛙与苹果的概率。

模型的第二部分则是通常的似然分布。它描述在每一个可能状态下应期待什么观测, 例如在是青蛙或是苹果的条件下, 会跳还是不会跳。把这两个部分结合起来, 有机体就能通过替身式地运行自己的生成模型, 对自己可能采取的行动或策略进行``如果......会怎样''的反事实模拟, 例如``如果我去池塘, 会发生什么?''。对状态做边缘化后, 我们就得到某项策略的边缘似然, 即该策略的证据 $P(\tilde{y} \mid \pi)$, 或者其自由能近似。换句话说, 只要知道策略如何影响状态转移, 我们就能够计算在该策略下某个观测序列发生的似然。正如公式 2.1 所示, 要计算某项策略的后验概率, 我们需要把这个似然与一个先验概率结合起来。

主动推断把这一规划问题分解为两个连续操作。第一步是给每个策略计算一个分数; 第二步是形成关于``应追求哪一策略''的后验信念。前者定义了关于策略的先验信念: 最好的策略拥有较高概率, 最差的策略则拥有较低概率。在主动推断中, 一项策略的优劣由其负期望自由能打分, 就像模型拟合的优劣由该模型的负自由能打分一样。策略的期望自由能 $G$ 与变分自由能 $F$ 不同, 因为计算前者必须考虑依赖于策略的未来观测; 相比之下, 后者只涉及当前与过去的观测。因此, 计算期望自由能时, 需要调用生成模型去预测在每个策略被执行时直至某个规划时域内将会出现的未来观测。并且, 因为一项策略会跨越多个时间步展开, 每项策略最终的期望自由能都必须对其所有未来时间步进行整合。

每项策略的期望自由能都可以通过取负号转换成一个质量分数, 并可直接作为主动推断主体的先验。这之所以成立, 是因为与物理学中的势能观念一致, 期望自由能是以对数概率空间来表达的。要把它转成关于策略的信念(或概率分布), 只需要进行指数化(以抵消对数)并归一化(以满足框 2.1 中的求和规则)。那些期望自由能更低的策略会被赋予更高概率, 从而成为有机体预期自己会采取的策略。

归根结底, 推断自己正在执行某项特定策略, 会对我们所预测的感觉数据产生后果。例如, 一个包含``屈肘''动作的策略, 会蕴含关于肱二头肌与肱三头肌本体感觉输入的预测。正是在这里, 规划与行动被连接起来: 与某个计划相关联的预测, 会转化为行动, 以解决它们与实际测得的本体感觉数据之间的差异(见第 2.3 节)。

\section{什么是期望自由能?}

到目前为止, 我们一直假定有机体在规划过程中会依照各策略的期望自由能来打分, 但我们还没有正面说明期望自由能究竟是什么。与变分自由能一样, 期望自由能也可以用多种在数学上等价的方式加以分解, 而每一种分解都为理解这一量提供了一个不同视角。

\begin{equation}
\begin{aligned}
G(\pi)
&= - E_{Q(\tilde{x}, \tilde{y}\mid \pi)} \Bigl[D_{\mathrm{KL}}\bigl[Q(\tilde{x}\mid \tilde{y}, \pi)\Vert Q(\tilde{x}\mid \pi)\bigr]\Bigr]
   - E_{Q(\tilde{y}\mid \pi)}[\ln P(\tilde{y}\mid C)] \\
&= E_{Q(\tilde{x}\mid \pi)}\bigl[H[P(\tilde{y}\mid \tilde{x})]\bigr]
   + D_{\mathrm{KL}}\bigl[Q(\tilde{y}\mid \pi)\Vert P(\tilde{y}\mid C)\bigr] \\
&\le E_{Q(\tilde{x}\mid \pi)}\bigl[H[P(\tilde{y}\mid \tilde{x})]\bigr]
   + D_{\mathrm{KL}}\bigl[Q(\tilde{x}\mid \pi)\Vert P(\tilde{x}\mid C)\bigr] \\
&= - E_{Q(\tilde{x}, \tilde{y}\mid \pi)}[\ln P(\tilde{y}, \tilde{x}\mid C)] - H[Q(\tilde{x}\mid \pi)]
\end{aligned}
\tag{2.6}
\end{equation}

这些分解中的第一种, 从直觉上说大概最有用, 因为它用与追求偏好观测(即利用, exploitation)完全相同的单位(nats), 来表达追求新信息(即探索, exploration)的价值。如此一来, 行为心理学中的经典``探索--利用''两难被消解了。通过最小化期望自由能, 这两项的相对权重就决定了行为是以探索为主, 还是以利用为主。需要注意, 语用价值(pragmatic value)在这里表现为关于观测的先验信念, 其中参数 $C$ 编码了偏好。先验信念与偏好之间这种可能显得不太直观的联系, 会在第 7 章中展开; 目前我们只需指出, 在把有价值的结果理解为每个主体所特有的那类结果(例如体温维持在 $37^\circ\mathrm{C}$)的前提下, 这一项可以被当作期望效用或价值。

信息增益项继承了我们在第 2.5 节讨论过的散度概念, 而正是这个散度保证了自由能是惊讶的上界。不过这里有一个转折: 我们现在不是要最小化散度, 而是要选择能最大化期望散度的策略, 也就是最大化信息增益。之所以会发生这种变化, 是因为我们现在是在对尚未观测到的未来结果取平均。这一点很微妙, 其核心在于结果的角色发生了切换。在评估观测结果的自由能时, 结果扮演的是后果; 但在评估期望自由能时, 结果则扮演原因的角色, 因为它们是未来尚未显现、却解释当前决策的隐藏变量。

由此得到的信息增益会惩罚那些从观测到状态是多对一映射的情形, 也就是说, 在不同状态下会产生相同观测, 从而无法实现精确信念更新。在人工智能与机器人学中, 能带来相同观测的状态(例如迷宫中两个看起来完全相同的 T 形岔路口)有时被称为别名化(aliased)状态, 而用简单方法(即只有刺激--反应, 没有推断或记忆)往往很难处理这种状态。问题在于, 仅凭当前观测, 我们无法知道自己到底处于哪个状态。主动推断则会由于这些状态的信息增益潜力较低, 而尽量避免把自己带入这样的局面。

一个简单例子有助于拆解信息增益(或认知价值)与语用价值之间的区别, 并说明为何在大多数现实情境中, 语用价值与认知价值必须被联合作为目标。想象一个想喝浓缩咖啡的人, 他知道镇上有两家不错的咖啡馆: 一家只在周一到周五营业, 另一家只在周末营业。如果他不知道今天是星期几, 那么他就必须先选择一个具有认知价值、能够消除不确定性的行动(例如看日历), 然后才能选择一个携带语用价值、能带来奖赏的行动(例如去正确的咖啡馆)。这个场景说明, 在大多数不确定情境中, 人们往往必须先执行认知性行动来消除不确定性, 然后才能有把握地选择语用性行动。那些不考虑选择所具认知可供性的策略选择方法, 只能通过随机数生成器来挑选策略, 结果往往就是错过那杯咖啡。因此, 只考虑语用价值的方案通常只能适用于不存在认知不确定性的情形, 例如这个人本来就知道今天是星期几, 因而可以直接前往正确的咖啡馆。

公式 2.6 的第二种分解把期望自由能写成风险(risk)与期望歧义(expected ambiguity)。这两项分别对应于复杂度与不准确性的类比: 风险就是期望复杂度, 歧义则是期望不准确性。风险是经济学中的常见概念, 对应的是策略与其后果之间可能存在的一对多映射, 即同一策略之下可能随机地产生几种不同结果。一个例子是带有随机奖赏的赌博场景(例如单臂老虎机), 其中你可能知道奖赏分布, 比如知道 $10\%$ 的时候会中奖。经济学把这称为风险情境, 因为在同样一次动作(拉动拉杆)之后, 你可能得到两种不同观测(中奖或未中奖)。这意味着, 你必须选择那些能够容纳不确定性的策略或计划。在像主动推断这样的风险敏感框架中, 关键是选择那些其概率性结果与先验偏好相匹配的策略, 这种匹配通过 KL 散度来衡量。简而言之, 当复杂度与偏离先验信念的程度是同一回事时, 最小化复杂度代价就等同于最小化风险。

同样地, 歧义对应于由于状态到结果之间映射不明确而导致的期望不准确性。若即便我们对生成结果的状态充满把握, 所预期的结果分布依然非常分散(或者说熵很高), 那么这一映射就是有歧义的。例如, 在知道天气是晴天还是雨天的条件下, 抛硬币出现正面还是反面的概率依然是 $50$--$50$, 因为天气与硬币结果之间没有关系, 因而这个映射具有最大歧义。于是, 仅仅看到一次反面并不能让我们获得关于天气的信息。需要注意的是, 大多数情境同时具有风险与歧义, 这意味着在状态到结果以及策略到结果之间都存在多对一或一对多的映射。回想一下, 结果(观测)是唯一真正可以被观测到的变量。主动推断能够自动处理这类情境, 因为期望自由能同时包含风险项与歧义项。

公式 2.6 的第三行则通过把风险重新表达为``关于状态的信念''与``以状态定义的偏好''之间的散度, 提供了另一种期望自由能的写法。这个形式的吸引人之处在于, 它可以像变分自由能(公式 2.5)那样被重排为期望能量与熵。不过, 这一写法的代价是它假定状态空间事先已知, 从而先验偏好可以直接与状态相联系。在大多数场景中, 这并不是问题, 而且用状态还是用结果来定义偏好, 在实践上差别也不大。不过, 更常见的做法还是以结果来定义偏好, 这样即使状态空间本身需要被学习, 外在动机也依然能够保留下来。

总之, 期望自由能既可以分解为风险与歧义, 也可以分解为语用价值与认知价值。这些分解很有意义, 因为它们使我们能够形式化地理解主动推断所处理的多种场景。此外, 它们也有助于我们看清主动推断如何吸收多种决策方案, 而这些方案通常可以通过忽略期望自由能中的一个或多个组成部分而得到(图 2.6)。如果移除先验偏好, 那么语用价值就变得无关紧要, 所有行动都只受认知可供性驱动, 此时这种方案只能处理不确定性的消解。在去掉先验偏好后, 负的期望自由能在不同语境里会被称为期望贝叶斯惊讶(例如注意性探索)或内在动机(例如自主学习)。如果移除歧义项, 得到的方案就对应于控制理论中的风险敏感控制或 KL 控制。最后, 如果同时移除歧义与先验偏好, 那么剩下的唯一要求就是最大化观测(或在第三行表述下的状态)的熵, 这可以被理解为不确定性抽样, 或者说``保持选择余地''。主动推断揭示了这些方案之间的形式关系, 也揭示了它们适用情境的局限。

尽管我们已经细致地分解了期望自由能, 以适应不同读者对这个泛函的理解方式, 但并不存在唯一正确或错误的切分方式。本书后半部分将会说明, 某一类自治系统只要要继续存在, 就必须选择那些看起来像是在最小化期望自由能的行动。从这个角度看, 认知性(探索性)命令与语用性(利用性)命令, 或风险与歧义, 都没有哪一方享有特权地位。这些可能人为制造出来的二分, 不过是同一枚存在性硬币的两面。

\section{低路的终点}

在引入了变分自由能与期望自由能这两个不同概念之后, 我们现在可以来看它们结合起来究竟实现了什么。这也标志着通往主动推断的低路走到了一个阶段性终点: 它从无意识推断这一观念出发, 经过贝叶斯大脑、知觉与行动的二重性, 最终到达``作为推断的规划''。

变分自由能是主动推断的核心。它衡量内部生成模型与当前和过去观测之间的拟合程度。通过最小化变分自由能, 生物体就最大化了自身模型的证据。这既保证生成模型会成为环境的良好模型, 也保证环境在一定程度上会服从该模型。

期望自由能则是一种给备选策略打分的方法, 用于规划。它在根本上是前瞻性的, 因为它考虑的是可能出现的未来观测; 同时它也是反事实性的, 因为这些可能的未来观测是以某个潜在可执行策略为条件的。期望自由能衡量的是行动策略相对于偏好中的未来状态与未来观测而言有多可信。通过按负期望自由能给策略打分, 参与主动推断的生物体实际上会相信自己正在追求那个使这一量最低的行动路线。从心理学角度说, 这意味着生物体关于策略的信念会直接对应于其意图, 而它会通过行动去实现这一意图。

从概念上看, 我们可以把变分自由能最小化与期望自由能最小化理解为两个推断回路, 其中一个嵌套在另一个内部。变分自由能最小化是主动推断的关键外层回路, 它足以优化知觉以及关于策略的信念。如果一个主动推断主体还拥有一个关于其行动后果的生成模型, 那么它就能够进一步评估期望自由能, 这便形成了内层回路。这种向未来进行规划的能力, 通过为策略赋予概率值, 支持了前瞻性的行动选择(Friston, Samothrakis, and Montague 2012; Pezzulo 2012)。

\section{总结}

主动推断是一种关于生命性人工物如何通过知觉与行动最小化惊讶, 或最小化惊讶的一个可处理代理量即变分自由能, 从而维系其存在的理论。在本章中, 我们试图从``知觉即推断''的贝叶斯处理方式出发, 再把这一思路扩展到行动领域, 以此为这一观念提供动机。贝叶斯推断依赖于一个关于感觉观测如何被生成的生成模型。这个模型以概率形式编码了有机体对世界的隐含知识, 具体表现为先验信念, 以及在不同状态与策略下所预期的结果。

主动推断的独特视角迫使我们重新审视贝叶斯推断中``先验''这一概念的通常语义。被期待的状态就是被偏好的状态, 其中包括支撑有机体生存的条件(例如特定生态位中的目标状态); 而与之相反, 令人惊讶的状态则是不被偏好的。通过实现自己的期待, 主动推断主体也就保障了自己的生存。正由于先验的概念与支撑有机体存在的条件之间存在紧密联系, 我们也可以说, 在主动推断中, 一个主体的身份与其先验是同构的。本书后面会让这种术语变得更加自然。

需要注意的是, 在这一视角下, 惊讶(surprise, 有时也写作 surprisal)是信息论中的一个形式构造, 并不必然等同于日常心理学意义上的``惊讶''。粗略地说, 有机体当前状态与先验(即编码其偏好状态的东西)偏离得越多, 该状态就越令人惊讶。因此, 主动推断的核心思想就是: 有机体(或它的大脑)必须主动最小化自己的惊讶, 才能维持生命。在某些条件下, 惊讶最小化可以被理解为减少模型与世界之间的差异。更一般地说, 主动推断真正最小化的量是变分自由能。变分自由能是惊讶的一个上界近似, 并且能够通过化学或神经层面的消息传递, 借助有机体生成模型可获得的信息被高效最小化。

重要的是, 知觉与行动以互补的方式共同最小化变分自由能: 前者通过改进自身的后验信念估计, 后者通过有选择地抽样那些被期望的观测。此外, 主动推断还会通过追随歧义与风险都最小的策略, 来最小化期望自由能。期望自由能因此把主动推断扩展到了前瞻性与反事实性的推断形式。至此, 我们沿着通往主动推断的低路完成了旅程。第 3 章中, 我们将转向高路: 它会从第一性原理与自组织出发, 抵达同样的结论。
